%!TEX root = main.tex

\chapter{Evaluation and Conclusion}
\label{ch:evaluation}

\section{Final evaluation against objectives}
\label{sec:eval-objectives}

Both objectives set out in Chapter~\ref{ch:introduction} were met, though not in the shape originally planned.

The first objective, preparing and using real multi-view and real low-light data, was met in full. The multi-view part went as planned. The lighting part did not, since it was planned as synthetic infrared-like augmentation and was replaced on purpose with training on real captured night footage once that data became available. This directly answers the interim supervisor feedback, which warned that synthetic low-light transforms should be described as an approximation of robustness, not a substitute for real infrared data. Rather than relabelling the experiment to satisfy this, it was replaced with one that does not need the caveat at all.

The second objective, a prototype on real, low-cost embedded hardware with an alert output reacting to sustained rather than momentary predictions plus a measured account of what getting there costs, was also met, though the hardware path changed along the way. The original plan treated a full FPGA-fabric CNN accelerator as the main goal, with a software fallback as backup. What was actually built runs the CNN on the ZC702's ARM core and only uses the fabric for the decision stage, a considered decision about scope, not a fallback settled for reluctantly. A dedicated accelerator is a multi-week job even starting from an existing architecture \citep{lu2025hardware}, and the IPR had already flagged this trade-off before it came up.

The quantification work answers that feedback most directly, since it asked that quantisation be checked through accuracy, model size, latency and throughput, before and after conversion, with laptop FPS clearly kept apart from hardware performance. Section~\ref{sec:qr-quantisation} reports all four measures with real numbers, and Table~\ref{tab:consolidated} gives the single combined comparison also asked for, with laptop and on-device rows kept separate rather than blended into one throughput claim.

\section{Project management}

The timeline set at the start (May to September 2026, single-view work through July, hardware work in August and September) fell short in one way, it did not expect robustness to become two full stages rather than one small check. Both were added after the qualitative domain-shift test in Section~\ref{sec:qr-baseline} made the problem real instead of theoretical. Looking back this was the right call, since the two robustness results are arguably the more original part of this project, but it did compress the hardware phase, a direct cause of later decisions, running the CNN on the ARM core instead of building a custom accelerator, and leaving the XNNPACK question unresolved. The lesson here is that a robustness finding worth investigating properly needs its own budgeted time from the start, not squeezed into an existing timeline once found. A second lesson is about ordering, since the lighting stage was the last big piece of work finished, leaving no room for it to go wrong. It did not, but that was not a margin the plan had earned.

\section{Insights gained}
\label{sec:eval-insights}

The most important technical insight was that a crash on unfamiliar embedded hardware is not necessarily proof the approach cannot work, it can just as easily be evidence of a specific, diagnosable, fixable mismatch between a generic published tool and the exact hardware in front of you. Giving up on on-device inference after the first \texttt{Illegal instruction} crash was a reasonable call given the information available at the time. Reopening that decision, working out the VFPv3-against-VFPv4 mismatch, and rebuilding the runtime from source on the board best shows the difference between treating an obstacle as a dead end and treating it as a problem with a findable cause.

The second insight is that robustness does not just show up on its own from seeing varied data. The frozen-backbone cross-view model saw both camera views during training and still showed a large, uneven gap, and the frozen-backbone low-light model saw both lighting conditions and still trailed by seven points on night data. Both problems needed the backbone itself to be adapted, and in both cases the real size of the problem would have stayed invisible if only the after-fix number had been measured. Reporting the ``before'' figure on purpose is what carries most of the analytical value here.

A third insight is that what looks like a model or hardware problem is sometimes a data-pipeline problem in disguise. Both the preprocessing mismatch (Section~\ref{sec:qr-runtime-correctness}) and the night-driver split bug were exactly this kind of error, and both would have produced confident, believable, wrong numbers instead of an obvious failure.

\section{Reflection on challenges}
\label{sec:eval-challenges}

Two more challenges, beyond the ARM32 runtime issue, are worth reflecting on honestly here, rather than just reporting how they were resolved.

The first is the WiFi adapter that was meant to give the board its own connectivity, and it was never made to work. The first idea was a USB power-budget limit, since several peripherals shared one port through unpowered hubs chained together, so it was tested directly. The competing devices were switched off, removing 1~A of demand, and the driver was re-bound with the whole bus available to it. The failure was byte-for-byte identical, the same four register-access timeouts at the same four offsets, real evidence against the power-budget theory, which shifted the diagnosis to a real driver or chipset incompatibility instead. The honest outcome is a disproven guess and a still-unfixed problem, deprioritised since on-device inference needs no network at runtime, with the value in having tested the guess rather than just asserting it.

The second is inference speed. The fix for the original crash, compiling with \texttt{-march=native}, should in theory also let XNNPACK's runtime CPU dispatch avoid the same mismatch, since it uses the same detection library already working correctly for the compiled NEON kernels, but ``should in theory'' is not the same as ``checked and confirmed''. Some XNNPACK ARM32 source files fix FPU flags directly rather than relying on runtime dispatch, and checking this properly would have needed a multi-hour rebuild that carried real risk to a working installation the demo depended on. Risking a working deployment for a speed gain this close to a deadline was not a good trade, so this is a reasoned decision not to chase the lead, not an unexamined gap.

\section{Comparison to the literature}

Chapter~\ref{ch:litreview} pointed out that the cross-view literature mostly reports post-fix accuracy with no equivalent unfixed baseline. This project's frozen-against-fine-tuned comparison (75.94\% to 91.76\%) points the same way as \citet{bhalla2025crossview} and \citet{sengar2023poseviNet}, even though both use more advanced methods. What this project adds is the size of the untreated gap on a standard architecture, something those papers do not report, and something needed to judge how much of any published improvement actually comes from the method rather than just from having multi-view data at all.

On lighting, the 19.53\% collapse is more extreme than the gradual drop \citet{wang2023driver} report across different recording conditions, and gives a real measured reference point for the concern behind \citet{shen2023visible}'s image-translation work. The problem their synthetic approach exists to reduce is severe enough to justify the effort, and the fact that real night data recovers 51.6 points suggests real data is still the stronger fix where it can be obtained. On hardware, \citet{lu2025hardware}'s 13.05 FPS for a dedicated accelerator on the same XC7Z020 chip puts this project's 0.11 FPS in perspective. That gap of roughly two orders of magnitude is the direct cost of the decision not to build one. This project's hardware contribution was never going to be a faster result, it is a numerically verified one instead.

\section{Limitations}
\label{sec:eval-limitations}

Run-to-run variance was not measured, since every result comes from a single training run with a fixed random seed, so no confidence intervals can be given and small differences between conditions should not be read as reliably meaningful. The interim feedback specifically asked for uncertainty across runs, and this report cannot give that honestly, since the data simply does not exist, and generating it would have meant re-running each stage several times, roughly two hours per cross-view run alone, which the remaining time did not allow. Stating this gap directly is the right response, rather than making up an interval or quietly ignoring the request. It matters most for the smaller comparisons. The 3.5-point day-accuracy cost and the 2.2-point margin between MobileNetV2 and EfficientNet-B0 both sit in a range where seed variance could plausibly matter, while the 51.6-point night recovery and 15.8-point cross-view improvement are far too large to be explained that way.

Test partitions are small in terms of subjects. Four held-out State Farm drivers and two low-light night drivers are still a small group, however many images they produce, since images from the same driver are not independent samples of anything.

Label quality in two datasets is not perfect, and this is documented rather than hidden. SAM-DD's class mapping was rebuilt by visually checking images, with a confidence level recorded for each row, and four low-light classes still carry an unresolved hand-side problem. No demographic information exists in any dataset used here, so no fairness claim can be made either way, and SAM-DD is recorded on a simulator, not in a moving vehicle. The FPGA results are post-synthesis estimates covering the alert stage only. Two further datasets, AUC V2 and 100-Driver, were requested but not granted within the project timeline, so cross-view robustness here rests on two sources rather than four.

\section{Commercial context and future work}
\label{sec:eval-future}

The commercial case sketched out in Chapter~\ref{ch:introduction} still holds, with one big qualification. Component costs stay realistic for a retrofit or fleet product, the alert logic takes up under 0.1\% of a low-cost FPGA, the quantised model is 2.6~MB, and the classifier runs on ordinary ARM hardware. But 0.11 FPS is not a workable product figure. At roughly nine seconds per inference and eight ticks in a row needed before an alert fires, this setup would take over a minute to respond to sustained distraction, and that gap has to close before any commercial claim is believable. The non-technical risks are, if anything, bigger, since driver monitoring feels like surveillance to a lot of people rather than help, and false alerts create liability risk while training users to ignore the system.

Four directions follow from this. First, inference speed. Turning XNNPACK back on is a clear next step, and beyond that Table~\ref{tab:consolidated} makes the case for a dedicated CNN accelerator, since the pipeline's whole cost sits in the classifier. Second, stronger cross-domain methods, using the contrastive viewpoint-invariant approach from \citet{bhalla2025crossview} on the unified dataset built here, with the same idea applying to the lighting domain too. Third, broadening the data, since the AUC V2 and 100-Driver requests would add two more independent camera sources, and the 45{,}600 daytime infrared images left out of the lighting experiment would allow a three-way comparison, day-RGB, day-infrared and night-infrared, separating the lighting gap from the sensor-modality gap. Fourth, repeated-seed runs, to supply the variance data this report has to report as missing.

\section{Conclusion}

This project set out to test three connected claims. The first was that a classifier's accuracy under one camera view says nothing reliable about its accuracy under another. The second was that the same is true, even more so, for lighting. The third was that a complete, hardware-verified detection-and-alert pipeline is actually achievable on genuinely low-cost embedded hardware, not just estimated for it. All three held up.

A MobileNetV2 model trained under subject-wise evaluation reached 84.67\% on State Farm's side view, but showed a real, uneven cross-view gap when simply pooled across views. Fine-tuning on real paired multi-view data closed most of that gap, up to 91.76\%. The same model collapsed to 19.53\% on real captured infrared-illuminated night footage, a failure severe enough that the original plan to approximate it synthetically would have understated it, and training on real night data recovered this to 71.13\%, at a measured cost of 3.5 points of daytime accuracy. Quantisation compressed the deployed model eightfold for 2.7 accuracy points, proved numerically identical on board and laptop, and cost roughly 150$\times$ in speed. And a hardware decision stage was verified against real model predictions, including real misclassifications, synthesised at almost no cost on a low-cost FPGA target.

What connects these results is not any single number, it is a methodological position. Robustness is worth measuring before applying the fix, on real data rather than approximations of it, and on the actual hardware the claim is about rather than whatever hardware is convenient. Each of those choices gave a less flattering figure than the alternative would have, and each gave a more defensible one. For a system whose failures would happen inside a moving vehicle, that seems like the right way round.
